% =====================================================================
%  Drop-in replacements for the red placeholder blocks in Dessert2026.
%  Every number here is produced by:
%     python3 gen_paper_tables.py --csv <results_FINAL.csv> --out paper/tables.tex
%     python3 gen_paper_figure.py --csv <results_FINAL.csv> --out paper/results_figure.pdf
%  Re-run both after any change to the CSV and these numbers stay in sync.
% =====================================================================


% ---------------------------------------------------------------------
% SECTION IV — IMPLEMENTATION.  Replaces "[model name(s), e.g., Claude
% Sonnet 4.5]" and "[container/VM sandbox, e.g., Docker]".
% ---------------------------------------------------------------------

Each of the four agents is realized as an LLM-backed module driven by
\texttt{claude-sonnet-5} at temperature 0.7, orchestrated as a LangGraph state
machine through a common controller that manages prompting, tool invocation,
and the per-agent execute-and-repair loop. Temperature is deliberately not
zero: the protocol calls for three runs per subject with different seeds so
that variance can be reported, and at temperature zero those runs would be
near-identical.

\textbf{Execution and sandboxing.} Local verification requires running
untrusted, machine-generated code. Each candidate clone is staged into an
isolated temporary directory and executed with the host interpreter under a
wall-clock timeout, capturing compiler diagnostics, exit status, standard
output, and assertion failures. We note this as a limitation rather than a
strength: the isolation is filesystem-level, not container-level, so a clone
that reached the network or the wider filesystem would not be contained. The
corpus is our own and the generated code is inspected, but a deployment on
untrusted subjects would require a container or VM boundary.

% ---------------------------------------------------------------------
% SECTION V.A — EVALUATION CORPUS.  Replaces "[Confirm the final program
% set and its provenance before submission.]"
% ---------------------------------------------------------------------

The corpus is twelve Python programs we authored for this study, modelled on
industrial measurement and control software: CRC framing, sensor
serialisation, timestamp parsing, SCADA polling, Modbus register access, MQTT
telemetry, PID control, a tank-level state machine, threshold alarming, and
three multi-module supervisory programs. Each ships a \texttt{pytest} suite and
recorded input/output traces that serve as both the \emph{VCA}'s test corpus
and the equivalence oracle's ground truth. Authoring rather than harvesting the
corpus was a deliberate trade: it costs external validity, and we treat that as
a threat in Section~VI, but it is what makes per-program golden vectors and
invariants available at all, and without those the equivalence question cannot
be asked at the granularity this study needs.

Two subjects, \texttt{modbus\_meter\_reader} (B2) and
\texttt{boiler\_control\_loop} (D1), declare a substitution that cannot be
satisfied as stated: both originals speak Modbus TCP via
\texttt{ModbusTcpClient(host, port=502)}, while the declared target
\texttt{minimalmodbus} is serial-only and exposes no TCP entry point. We
discovered this during evaluation rather than by design. Rather than repair the
corpus we retained both as \emph{negative controls}, and they proved
informative: across the nine D1 runs, five independently chose
\texttt{pyModbusTCP} over the declared target, and the single run that did
adopt \texttt{minimalmodbus} was caught by the factory-execution check passing
a TCP port number as a Modbus slave address.

% ---------------------------------------------------------------------
% SECTION V.B — METRICS.  Replaces the structural-divergence definition.
% ---------------------------------------------------------------------

\textbf{Structural divergence} is reported as a composite of two
rename-insensitive measures taken between the delivered clone and the pristine
original, over the concatenation of all first-party files. The first is the
Ratcliff/Obershelp similarity of the pre-order AST \emph{node-type} sequence,
which ignores identifier and literal choices and so cannot be inflated by
renaming. The second is a weighted Jaccard distance over a depth-tagged
histogram of branch constructs, which captures reshaped control flow. The
composite is their unweighted mean, on $[0,1]$, higher meaning more divergent.

We report a rename-insensitive measure deliberately. Token-level edit distance
rewards renaming every identifier, which is exactly the kind of cosmetic change
a Type~II clone detector defeats, and a diversification result that a detector
defeats is not a result.

% ---------------------------------------------------------------------
% SECTION V.C — EXPERIMENTAL SETUP.  Replaces the model/temperature/caps,
% the requirement-set description, and the two baseline descriptions.
% ---------------------------------------------------------------------

We run \emph{Proteus} with \texttt{claude-sonnet-5} at temperature 0.7, three
seeds per subject, a wall-clock ceiling of 180\,s per LLM request, and a cap of
three \emph{VCA}-triggered stage re-entries per run. Every subject receives the
same requirement: remove the dependency named in its \texttt{requirements.txt}
while preserving observable behaviour, and additionally, for Categories C and
D, diversify all first-party logic.

Crucially, \textbf{the replacement library is never named} --- not to
\emph{Proteus} and not to the baseline. Both must select a target themselves.
An earlier protocol supplied the corpus's declared target to \emph{Proteus}
while the baseline chose freely; that made the comparison asymmetric and, on
the two negative controls, forced \emph{Proteus} toward a library that cannot
work. Under free choice the two systems are asked the same question.

\textbf{Baselines.} To isolate the contribution of \emph{Proteus}'s
verification-driven design we compare against two configurations.
\emph{Baseline~1 (single-shot LLM)} is one prompt to the same model, naming
only the dependency to be removed, with no agents, no tools, and no
verification: a single API call whose reply is parsed for code blocks. We
confirmed by inspection and by telemetry that it uses no tool loop --- exactly
one LLM call on all 36 runs, averaging 1{,}471 input tokens, against up to 14
calls and 20{,}749 input tokens per run for \emph{Proteus}. \emph{Baseline~2
(\emph{Proteus} without verification)} is an ablation with the \emph{VCA} and
the global oracle disabled, leaving only the per-agent local execute-and-repair
loop.

% ---------------------------------------------------------------------
% SECTION V.D — RESULTS.  Replaces the whole red placeholder block.
% ---------------------------------------------------------------------

Table~\ref{tab:results} reports the conventional measures. They do not separate
the three configurations: generation success is 30/36, 33/36 and 29/36, and
equivalence confirmation on Categories C and D is 16/18, 17/18 and 16/18. On
these numbers alone the ablation is the strongest configuration and full
\emph{Proteus} the weakest, which is the opposite of what the two-tier
discipline predicts.

The measures are not wrong; they are answering a narrower question than they
appear to. Every subject's tests pass before migration, and every subject
injects its I/O dependency, so the factory code that performs the substitution
is frequently never executed by the suite. A clone can therefore compile, pass
its full suite, and satisfy the oracle while calling the new library
incorrectly on the one path the tests do not reach.

To measure what those instruments cannot, we submitted all 108 clones to a
blind LLM review (Table~\ref{tab:judge}). Each clone was shown alongside its
original to \texttt{claude-opus-5} --- deliberately a different and stronger
model than the \texttt{claude-sonnet-5} that generated it, so that no model
grades its own output --- and asked only for concrete behavioural differences,
each naming a file, quoting a line, and stating the runtime consequence.
Verdicts are \textsc{broken} (fails on ordinary input), \textsc{major} (wrong
behaviour in a case the original handled), or \textsc{minor} (a real difference
in an edge case). All system-identifying markers were renamed or removed before
review, and the batch aborts rather than continuing if any marker survives, so
the reviewer never learns which configuration produced a clone.

The review separates the configurations where the conventional measures do
not. Of the clones each configuration \emph{itself accepted} ---
\texttt{migration\_success}: compiles, passes its suite, genuinely changed, and
the dependency verifiably removed --- the fraction carrying \textsc{broken} or
\textsc{major} defects is 54\% for the single-shot baseline, 55\% for the
ablation, and 33\% for full \emph{Proteus} (Table~\ref{tab:precision},
Fig.~\ref{fig:results}).

Two observations follow, and they point in opposite directions.

First, \textbf{the ablation performs like the single-shot baseline}: 55\%
against 54\%. Removing the verification tier leaves an agent architecture that
is no more reliable than one prompt. Whatever \emph{Proteus} contributes is
attributable to verification rather than to decomposition into agents.

Second, \textbf{full \emph{Proteus} does not deliver more sound clones than one
prompt}. It yields 10 defect-free accepted clones per 36 runs against the
baseline's 12. It accepts far fewer clones (15 against 28) and is right more
often when it does accept, which is a precision gain paid for in recall. On a
corpus where a wrong clone is worse than no clone --- boiler supervisors, pump
controllers --- that is the trade we would choose, but it is a trade and not a
dominance result.

\textbf{Cost.} The precision is not free. Full \emph{Proteus} consumed 925K
tokens and 140 minutes across 36 runs against the baseline's 95K tokens and
6 minutes, roughly a sixfold increase in tokens and a twentyfold increase in
wall time, for \$4.91 against \$0.79.

\textbf{Statistical caution.} With twelve subjects the differences above are
not statistically significant: Fisher's exact test gives $p=0.34$ for the
single-shot baseline against full \emph{Proteus} on defect rate, and $p=1.00$
against the ablation. The 95\% confidence intervals overlap substantially
(54\%~[36--70], 55\%~[28--79], 33\%~[15--58]). We report the direction and its
mechanism, not a demonstrated effect; separating 33\% from 54\% at $p<0.05$
would require roughly three to four times this corpus.

% ---------------------------------------------------------------------
% SECTION V.E — ANALYSIS.  Replaces the red placeholder block.
% ---------------------------------------------------------------------

\textbf{Which stage dominates repair cost.} Full \emph{Proteus} triggered six
\emph{VCA} stage re-entries across 36 runs, all routed to the \emph{GSA}. No
run exhausted its repair budget. Repair cost is therefore not the binding
constraint on this corpus; acceptance is.

\textbf{Which failure classes are hardest.} The blind review's findings fall
into a small number of recurring shapes, and the two systems fail differently.
The single-shot baseline produces 80 findings dominated by \emph{fidelity}
failures: 34 concern a changed exception type or message, and 10 concern a
client or connection that is never closed. Its wiring is correct and its
semantics drift. \emph{Proteus} produces 45 findings dominated by
\emph{integration} failures: arguments accepted by an adapter and then never
used (7 against the baseline's 2), adapters that are generated but not reached
(5 against 2), and errors collapsed into a single generic exception (4 against
1). Its adapters are structurally sound --- generating a wrapper naturally
preserves lifecycle and error translation, which is why it has one unclosed
resource against the baseline's ten --- and it loses detail at the call sites.

\textbf{Two illustrative defects.} In
\texttt{proteus\_no\_verif/D1\_boiler\_control\_loop\_run2/transport.py}, an
adapter catches \texttt{(ModbusError, Exception)} and converts every failure
--- connection loss, and programming errors such as \texttt{TypeError} alike
--- into a generic \texttt{IOError}. On a boiler supervisor an adapter bug
becomes indistinguishable from a genuine Modbus fault. The ablation accepted
this clone; full \emph{Proteus} rejected its counterpart at the \emph{VCA} on
the same two seeds. In
\texttt{proteus\_full/D1\_boiler\_control\_loop\_run3/transport.py}, the
generated wrappers accept the per-call slave identifier and never use it, so an
injected client configured for a different device would be actuated instead.
Neither defect is reachable by the test suite, because the suite injects a fake
client that ignores the same argument.

\textbf{What the global oracle catches.} Two \emph{Proteus} clones import the
migrated-to library while still calling the migrated-from one, which was never
imported --- a guaranteed \texttt{NameError} on the factory path. Both were
rejected by the pipeline, at the \emph{LMA} and \emph{VCA} respectively, and
the blind reviewer independently rated both \textsc{broken}. The system's own
verdict and an independent reviewer's agree on exactly the clones the test
suite passes.

\emph{a) Lessons Learned:} Three findings we did not anticipate.

\emph{Test suites are the wrong instrument for this task.} Dependency injection
is good practice and it makes the migrated code path untestable by the very
suite that certifies it. Every configuration compiles at 36/36; compilation
discriminates nothing. Any future benchmark of clone generation needs an
oracle that executes the substituted path directly, which is why we added
factory execution as a fifth layer mid-study.

\emph{Separating substitution from mutation keeps repairs local, but the seam
leaks.} The \emph{LSA}/\emph{GSA} split behaves as intended --- repairs stay
within one agent and no run exhausted its budget --- yet the residual failure
mode lives precisely at the boundary, where an adapter is generated correctly
and then not fully wired to the call sites it was written for. Orthogonal
concerns still require a shared contract at their interface.

\emph{One third of the corpus is unsolved by every configuration.} Four
subjects --- the MQTT bridge, the PID controller, the threshold engine, and the
telemetry gateway --- produced zero defect-free accepted clones across all 27
runs. Category C is 1 of 27 across all configurations. The requirement to
replace a dependency with first-party code is, on this evidence, not
differentially solved but essentially unsolved: \emph{Proteus}'s resolver
selects among libraries and has no mechanism to synthesise a replacement
implementation, and while the single-shot baseline does write a PID controller
and a state machine by hand, only one such clone survived review.

% ---------------------------------------------------------------------
% ABSTRACT and CONCLUSIONS — the [X] placeholders.
% ---------------------------------------------------------------------

% ABSTRACT, final sentences:
\emph{Proteus} produces working clones for 29 of 36 runs across twelve subject
programs at a median of zero repair iterations, and the oracle confirms
behavioural equivalence on 16 of 18 Category~C and~D runs. A blind review of
all 108 generated clones finds that 33\% of the clones \emph{Proteus} accepts
carry serious defects, against 54\% for an unverified single-shot baseline ---
though \emph{Proteus} accepts far fewer clones, and on this corpus size the
difference is not statistically significant. These early results suggest that
agentic, verification-driven mutation is a promising route to automated,
trustworthy software diversification, and that the verification tier rather
than the agent decomposition is what carries the effect.

% CONCLUSIONS, the [X] sentence:
On software that estimates factory measurements and controls industrial
infrastructure, \emph{Proteus} generates equivalence-confirmed clones for 16 of
18 Category~C and~D runs at a median of zero repair iterations, and reduces the
share of accepted clones carrying serious behavioural defects from 54\% to
33\% relative to an unverified single-shot baseline --- a difference in the
predicted direction that this corpus is too small to establish
($p=0.34$).
